\paragraph{Follow purchasing, sales and wages through one book.}
The next declared methods connect six ordinary operations: purchasing inventory
on credit, delivering it on credit, collecting receivables, paying suppliers,
accruing wages and paying those wages. Each method is a data definition interpreted
by the same checked requirements, monetary and memory engines. An explicit role
map identifies the accounts. Quantity and unit cost carry different dimensions,
so their product is a currency amount; substituting a money input for units fails
the type check. The selected synthetic premise is a positive, constant inventory
unit cost stored in company memory. It does not establish an inventory valuation
policy for heterogeneous lots.

Purchasing increases inventory and supplier payables without moving cash.
Delivery increases receivables and revenue, while transferring its cost from
inventory to cost of sales. Collection reduces receivables and increases cash.
Supplier and wage payments reduce their respective liabilities and cash; wage
accrual creates the expense and payable before payment. The methods check available
inventory, receivables, obligations and cash at the event's actual turn. They also
retain an explicit applicability finding: unknown remains unresolved, and false
remains unsupported. Other accounts and all typed memory remain unchanged.

Start with cash USD~100, inventory USD~20 and equity USD~120, at a unit cost of
USD~5. Purchasing six units creates inventory and payables of USD~30. Delivering
five units for USD~50 creates that receivable and releases USD~25 of inventory.
Collect USD~35, pay the supplier USD~20, accrue wages of USD~10 and pay USD~8.
Cash ends at USD~107, receivables at USD~15 and inventory at USD~25. Supplier
payables of USD~10 and wage payables of USD~2 leave equity of USD~135, including
profit of USD~15. Assets of USD~147 equal liabilities of USD~12 plus equity.
Operating cash of USD~7 also reconciles to profit less the USD~15 receivable
increase and USD~5 inventory increase, plus the two payable increases of USD~10
and USD~2. This bridge explains the supplied cycle; it does not select a general
indirect cash-flow policy.

Every journal, complete account balance and cash movement agrees with separate
calls to the retained \texttt{DemoCorporate/1} methods under the same premises.
A second chart produces identical financial quantities through its own explicit
mapping, and the text frontend reproduces the data program. Additional checks
retain nonzero plant, accumulated depreciation and debt while these operations
run. This is monetary conformance for the selected methods, not a conversion of
all native asset, history, tax or information-selection semantics.

A payment exceeding available cash preserves the preceding book. The current
requirement interpreter returns a primitive invalid result for its failed cash
condition. A separate diagnostic identifies an isolated liquidity shortage as
an economic funding failure and retains that exact underlying result; another
failed condition is never hidden by that explanation. A general typed default
outcome remains a separate execution obligation. Maturity, ageing, impairment,
tax and disclosure completeness also remain unresolved until their own facts and
methods are supplied. The next step is to derive explicit working-capital
roll-forwards from these verified events, rather than treating balanced totals
as complete notes.

